\cleardoublepage
\addcontentsline{toc}{chapter}{Bibliografia}

Adamski J., Czarnecki R.: Tytuł artykułu, Czasopismo, 1/2004, s. 21-23.

Bartecki K.: Classical vs. neural network-based PCA approaches for lossy image compression: Similarities and differences. Applied Soft Computing, vol. 161, No. 111721, 2024.

Morecki A., Knapczyk J. (red.): Podstawy robotyki. WNT, Warszawa, 1993.

Przelaskowski A.: Kompresja danych, skrypt internetowy: http://www.ire.pw.edu.pl. Dostęp: 10 października 2024.

Sciavicco L., Siciliano B.: A Solution Algorithm to the Inverse Kinematic Problem for Redundant Manipulators. IEEE Journal for Robotics and Automation, vol. 4, issue 4, pp. 403-410, 1988.

Skarbek W.: Multimedia. Algorytmy i standardy kompresji, Akademicka Oficyna Wydawnicza PLJ, Warszawa, 1998.

Typowy wykaz literatury będzie liczył od kilku do kilkunastu pozycji, posortowanych alfabetycznie według nazwisk autorów. Należy dokładnie zastosować pokazany wyżej sposób jego formatowania. Należy zachować szczególną ostrożność przy korzystaniu z artykułów zamieszczonych w Wikipedii jako źródła wiedzy dotyczącej specjalistycznych dziedzin, np. automatyki czy robotyki. Niekiedy są one niskiej jakości, a nawet zawierają błędy. Bardziej wartościowe wersje haseł często można znaleźć np. w angielskojęzycznych odpowiednikach jej artykułów. Powyższa uwaga dotyczy również korzystania z różnego rodzaju chatbotów, wykorzystujących tzw. duże modele językowe oraz metody uczenia maszynowego. Najlepszym rozwiązaniem jest odwoływanie się do materiałów źródłowych w postaci wysokiej jakości opracowań, recenzowanych przez specjalistów z danej dziedziny: np. z książek, artykułów naukowych i artykułów z czasopism (również dostępnych w wersjach cyfrowych).

% \nocite{*}
% \bibliographystyle{plain}
% \bibliography{bibliografia}

\printbibliography